% Part:sets-functions-relations
% Chapter: size-of-sets
% Section: introduction

\documentclass[../../../include/open-logic-section]{subfiles}

\begin{document}

\olfileid[pt]{sfr}{siz}{int}

\olsection{Introdução}

Quando Georg Cantor desenvolveu a teoria dos conjuntos na década de 1870, um dos
seus objetivos era tornar aceitável a ideia de uma coleção infinita---um
infinito \emph{atual}, como diriam os medievais. Parte central disto foi o
seu tratamento do \emph{tamanho} de conjuntos diferentes. Se $a$, $b$ e $c$ são todos distintos, então o conjunto $\{a, b, c\}$ é intuitivamente
\emph{maior} que $\{a, b\}$. Mas e os conjuntos infinitos? Serão todos do mesmo
tamanho? Acontece que não.

A primeira ideia importante aqui é a de \emph{enumeração}. Podemos
listar qualquer conjunto finito enumerando todos os seus !!{element}s. Para
alguns conjuntos infinitos, também podemos listar todos os seus !!{element}s se
permitirmos que a própria lista seja infinita. Tais conjuntos chamam-se
!!{enumerable}s. O resultado surpreendente de Cantor, que compreenderemos
completamente no fim deste capítulo, foi que alguns conjuntos infinitos \emph{não} são
!!{enumerable}s.

\end{document}
